% !TeX root = ../../Book1.tex
\chapter{高阶几何测度论与流}
\label{b1:app:D}

\begin{chapterabstract}
可整流集合告诉我们，一个粗糙集合在几乎每点仍可以由切平面描述；
BV 理论进一步把集合的边界写成向量测度。
本附录加入取向和整数重数，将这些对象放进一个可以取边界、推前和弱极限的线性空间。
目标不是罗列流理论的全部定理，而是从微分形式的积分出发，
理解为什么质量与边界质量共同受控时，可以取得保留整数几何结构的极限，
并用它建立几何变分问题的存在性。
\end{chapterabstract}

先修材料包括\chapref{b1:ch:19}的可整流集及面积、余面积公式，
\chapref{b1:ch:20}的分布，以及\chapref{b1:ch:27}的有限周长集合。
附录\ref{b1:app:B}提供外代数、Jacobi 因子与取向约定。
切测度部分只需 Radon 测度的微分与淡收敛；
从流的定义开始，则要把测试函数换成交错微分形式。

\begin{learninggoals}
\begin{enumerate}
 \item\label{b1:item:app-d-goal-tangents}
 区分测度放大的尺度系数与映射推前，理解可整流测度的平面切测度。
 \item\label{b1:item:app-d-goal-objects}
 能从测试形式定义流、边界和质量，分清正规、可整流与整数重数三个条件。
 \item\label{b1:item:app-d-goal-operations}
 掌握推前、同伦、切片的边界公式，并能用简单曲面核对取向符号。
 \item\label{b1:item:app-d-goal-compactness}
 分清正规流的解析抽取与整数闭合的几何内容，
 用平坦范数、变形和积分流紧性处理填充及最小化问题。
\end{enumerate}
\end{learninggoals}

主要伴读是 Simon 的《Lectures on Geometric Measure Theory》%
\cite{Simon1983Lectures}第26--32节。
其中微分形式上的泛函、切片、网格变形和整数闭合构成连贯路线。
本附录完整证明基本构造、正规流切片估计、解析紧性以及变分应用；
变形定理和整数闭合定理的完整结构证明较长，给出明确出处及证明概要。
后两项是采用的深层输入，不能把对它们的引用误当成本附录已独立重建整套流理论。

% !TeX root = ../../Book1.tex
\section{切测度与局部放大}
\label{b1:sec:D101}

研究一张光滑曲面在一点附近的形状，可以先平移该点到原点，再把距离放大。
在合适的面积归一化下，曲面测度趋于切平面上的面积测度。
对于尚不知道是否可整流的测度，不能预先指定极限一定是平面；
应当先容许一般非零测度作为放大极限，再研究这些极限的结构。
这一区别使切测度既能描述已有的切平面，也能用于寻找切平面。

\subsection{归一化与非零极限}
\label{b1:subsec:D10101}

设 \(\mu\) 是 \(\mathbb R^n\) 上的正 Radon 测度，\(x\in\operatorname{supp}\mu\)。
对 \(r>0\)，记
\[
 \eta_{x,r}(y)=\frac{y-x}{r}.
\]
这里的推前不带 Jacobi 因子：
\[
 (\eta_{x,r})_\#\mu(A)=\mu(x+rA).
\]
尺度归一化另由一个正数给出，不能在推前的定义和前置系数中重复计算。

\begin{definition}\label{b1:def:app-d-tangent-measure}
若存在 \(r_j\downarrow0\) 和 \(c_j>0\)，使某个非零正 Radon 测度 \(\nu\) 满足
\[
 c_j(\eta_{x,r_j})_\#\mu\longrightarrow\nu
 \quad\text{在淡收敛意义下},
\]
则称 \(\nu\) 为 \(\mu\) 在 \(x\) 处的一个%
\term[qiecedu]{切测度}[tangent measure]。
所有这些测度组成的集合记为 \(\operatorname{Tan}(\mu,x)\)。
\symidx{\(\operatorname{Tan}(\mu,x)\)：测度在一点的切测度族}
\end{definition}

淡收敛由所有连续紧支集函数测试，沿用%
\cref{b1:def:vague-convergence}。
一般的切测度不是概率测度，定义中的 \(c_j\) 也不必等于某个预先指定的 \(r_j^{-k}\)。
排除零测度很重要：把 \(c_j\) 取得足够小，总能制造出没有几何信息的零极限。
另一方面，归一化不能自动保证紧性；放大后的质量还须在每个固定紧集上一致受控。

例如，若 \(\mu=\theta\mathcal H^k\llcorner P\)，其中 \(P\) 是经过 \(x\) 的 \(k\) 维仿射平面且 \(\theta>0\)，则
\[
 r^{-k}(\eta_{x,r})_\#\mu
   =\theta\mathcal H^k\llcorner(P-x).
\]
若 \(\mu\) 在 \(x\) 处有原子，固定系数放大产生的则是点质量，而非正维数的平面面积。
两个例子说明，尺度幂反映了测度在该点的维数，不能只根据环境空间的维数决定。

\begin{proposition}\label{b1:prop:app-d-tangent-scaling}
若 \(\nu\in\operatorname{Tan}(\mu,x)\)，则对任意 \(a,s>0\)，
\(a(\eta_{0,s})_\#\nu\) 也属于 \(\operatorname{Tan}(\mu,x)\)。
\end{proposition}
\begin{proof}
取定义中的 \(c_j,r_j\)。对 \(\varphi\in C_c(\mathbb R^n)\)，
\(\varphi\circ\eta_{0,s}\) 仍是连续紧支集函数，而且
\[
 \eta_{0,s}\circ\eta_{x,r_j}=\eta_{x,sr_j}.
\]
于是
\[
 ac_j(\eta_{x,sr_j})_\#\mu
   \longrightarrow a(\eta_{0,s})_\#\nu.
\]
极限非零，因为伸缩是同胚且 \(a>0\)。
\end{proof}

这条封闭性不表示每个切测度本身都齐次，也不表示切测度族对加法封闭。
“一个测度的所有伸缩仍在该族内”与“该测度在伸缩下只乘一个幂次”是不同断言。

\subsection{密度控制如何产生切测度}
\label{b1:subsec:D10102}

一般切测度理论需要细致选择尺度。本节先证明一个可直接使用的充分条件，
它已覆盖正有限密度点，也说明归一化和紧性各自承担什么工作。

\begin{proposition}\label{b1:prop:app-d-tangent-existence-density}
设 \(k>0\)。若存在 \(a,b,r_0>0\)，使
\[
 ar^k\leq\mu(B(x,r))\leq br^k
 \quad(0<r<r_0),
\]
则任意趋于零的尺度序列都有子列，使
\(r_j^{-k}(\eta_{x,r_j})_\#\mu\) 淡收敛于一个非零 Radon 测度。
\end{proposition}
\begin{proof}
记这些归一化测度为 \(\nu_j\)。
对每个固定 \(R>0\)，当 \(j\) 足够大时，
\[
 \nu_j(B(0,R))=r_j^{-k}\mu(B(x,Rr_j))\leq bR^k.
\]
有限个较早指标各自在这个球上也有有限质量。
因此 \((\nu_j)\) 在每个紧集上质量一致有界。
在依次增大的闭球上，对 \(C_c(\mathbb R^n)\) 的可数一致稠密测试族作对角抽取，
所得极限是局部有界的正线性泛函；由 Riesz 表示得到正 Radon 测度 \(\nu\)。
局部质量界把稠密测试族上的收敛推广到全部连续紧支集函数。

再取 \(0\leq\varphi\leq1\)，使 \(\varphi=1\) 于 \(\overline{B(0,1)}\)，
且 \(\operatorname{supp}\varphi\subset B(0,2)\)。
充分大的 \(j\) 满足 \(\int\varphi\,d\nu_j\geq a\)，
故 \(\int\varphi\,d\nu\geq a>0\)。
这一步使用下密度界，保证极限不因质量消失而退化为零。
\end{proof}

上密度界负责阻止固定观察窗口内质量无限增大，下密度界负责保留非零对象。
如果只写“把球的质量归一化为一”，仍须检查较大球上的质量；
一个单位球中的概率归一化并不控制所有固定半径的球。
本命题不声称是切测度存在的必要条件，也不替代一般 Radon 测度几乎处处的切测度理论。

\subsection{可整流测度的平面极限}
\label{b1:subsec:D10103}

\chapref{b1:ch:19}已给出可整流集和带权测度的密度结论。
把这些结论写成切测度语言，可以看清密度、权重和切平面如何同时进入极限。

\begin{theorem}\label{b1:thm:app-d-rectifiable-tangents}
设 \(1\leq k\leq n\)，\(E\subset\mathbb R^n\) 可数 \(k\) 维可整流，
\[
 \mu=\theta\mathcal H^k\llcorner E
\]
是 Radon 测度，其中 \(\theta\) 可测、非负。
则对 \(\mu\)-几乎每个 \(x\)，存在 \(k\) 维线性空间 \(P_x\)，使
\[
 r^{-k}(\eta_{x,r})_\#\mu
   \longrightarrow\theta(x)\mathcal H^k\llcorner P_x
 \quad(r\downarrow0),
\]
且
\[
 \operatorname{Tan}(\mu,x)
   =\{\,a\mathcal H^k\llcorner P_x:a>0\,\}.
\]
这里选取 \(\theta\) 的有限正值代表；结论只涉及满足该条件的满 \(\mu\) 测度集合。
\end{theorem}
\begin{proof}
先在 \(E\) 局部面积有限的情形工作。
由\cref{b1:thm:rectifiable-density-tangent}，在几乎每个 \(x\in E\)，
无权面积的归一化放大趋于 \(\mathcal H^k\llcorner P_x\)。
同时，由 Radon 测度的微分定理，几乎每个这样的点满足
\[
 \int_{E\cap B(x,r)}|\theta(y)-\theta(x)|\,d\mathcal H^k(y)
     =o(r^k).
\]
对支集包含于 \(B(0,R)\) 的测试函数 \(\varphi\)，有
\[
 \begin{split}
 &\left|r^{-k}\int_E
   \varphi\!\left(\frac{y-x}{r}\right)
   \bigl(\theta(y)-\theta(x)\bigr)\,d\mathcal H^k(y)\right|\\
 &\hspace{1em}\leq
  \|\varphi\|_\infty\,r^{-k}
  \int_{E\cap B(x,Rr)}|\theta(y)-\theta(x)|\,d\mathcal H^k(y)
  \longrightarrow0.
 \end{split}
\]
由无权放大极限便得到第一式。

一般的 \(E\) 未必局部面积有限。
如\cref{b1:prop:rectifiable-weight-density}的证明，
把它分解为可数个互不相交、有限面积的 Borel 片 \(E_\ell\)，
再按位置和权重截断。
对 \(\mu\)-几乎每个 \(x\in E_\ell\)，测度微分定理给出
\[
 \frac{\mu(B(x,r)\setminus E_\ell)}{\mu(B(x,r))}\longrightarrow0,
\]
而同一处带权密度满足
\(\mu(B(x,r))/(\omega_kr^k)\longrightarrow\theta(x)\in(0,\infty)\)。
因而片外部分为 \(o(r^k)\)；
在包含测试函数支集的 \(B(x,Rr)\) 上使用同一估计，
即可把片上的极限推广到整个 \(\mu\)。

最后设 \(c_j(\eta_{x,r_j})_\#\mu\to\nu\neq0\)，并令
\(\lambda=\theta(x)\mathcal H^k\llcorner P_x\)。
已知 \(r_j^{-k}(\eta_{x,r_j})_\#\mu\to\lambda\)。
取非负测试函数 \(\varphi\)，使 \(\int\varphi\,d\lambda>0\)。
若 \(c_jr_j^k\) 无界，则对应子列在这个测试函数上的积分无界，与淡收敛矛盾。
若它有趋零子列，则每个测试函数的积分沿该子列趋零，与 \(\nu\neq0\) 矛盾。
故可再取子列，使 \(c_jr_j^k\to c\in(0,\infty)\)，于是 \(\nu=c\lambda\)。
反向包含由第一式乘以任意正常数得到。
\end{proof}

这是一条几乎处处的结论。
两条相交直线上的长度测度在交点放大后仍是两条直线的长度测度之和；
它不是单一平面测度，但交点对长度测度为零，不违反定理。
反过来，由正有限密度推出可整流性，需要%
\cref{b1:thm:preiss-rectifiability}中的深层结果。
本节只把已经可整流的对象的局部放大写清楚，不把这个反向过程当成紧性抽取的直接后果。

切测度只记录质量分布，没有记录曲面的取向。
即使知道某点的全部切测度，也不能从中恢复曲面应当朝哪个方向积分。
接下来引入的流把这种取向信息与边界运算放在同一个线性框架中。

% !TeX root = ../../Book1.tex
\section{微分形式、流与质量}
\label{b1:sec:D102}

对一条有向曲线积分时，反向参数化改变积分的符号，却不改变曲线的长度。
高维曲面也有同样的两种信息：积分带符号，面积不带符号。
流首先保存前者，再通过对所有单位测试形式取上确界恢复后者。
因此它与正测度的差别不只在于允许负数，还在于测试对象携带维数和方向。

\subsection{测试形式与外微分}
\label{b1:subsec:D10201}

本节在开集 \(U\subset\mathbb R^n\) 中工作。
\appref{b1:app:B}已构造外空间及其对偶配对。
记 \(I=(i_1,\ldots,i_k)\) 为严格递增多重指标，并置
\(dx^I=dx^{i_1}\wedge\cdots\wedge dx^{i_k}\)。

\begin{definition}\label{b1:def:app-d-differential-form}
若
\[
 \omega(x)=\sum_{|I|=k}\omega_I(x)\,dx^I,
 \quad \omega_I\in C^\infty(U),
\]
则称 \(\omega\) 为 \(U\) 上的光滑 \(k\) 次%
\term[weifen-xingshi]{微分形式}[differential form]。
所有具有紧支集的此类形式组成 \(\mathcal D^k(U)\)。
其收敛约定为：支集包含在同一个紧集内，且每个系数的每阶导数都一致收敛。
\end{definition}

零次形式就是函数，零次外空间约定为 \(\mathbb R\)。
外积的归一化沿用\cref{b1:prop:app-b-form-pairing}，不额外乘 \(1/k!\)。
定义%
\term[wai-weifen]{外微分}[exterior derivative]为
\[
 d\omega=\sum_{|I|=k}\sum_{i=1}^n
       \partial_i\omega_I\,dx^i\wedge dx^I.
\]
对最高次形式规定 \(d\omega=0\)。

\begin{proposition}\label{b1:prop:app-d-exterior-derivative}
外微分是测试形式空间之间的连续线性映射，并满足
\[
 d^2\omega=0,\quad
 d(\alpha\wedge\beta)
   =d\alpha\wedge\beta+(-1)^p\alpha\wedge d\beta
\]
其中 \(\alpha\) 为 \(p\) 次形式。
\end{proposition}
\begin{proof}
连续性来自系数的一阶微分：输出的每个固定阶导数由输入的下一阶导数控制，
支集不扩大。
计算 \(d^2\omega\) 时，\(i=j\) 的项含 \(dx^i\wedge dx^i=0\)；
当 \(i\neq j\) 时，交换 \(i,j\) 不改变二阶偏导，却使外积改变符号，所以两项抵消。
对第二式只需考察
\(\alpha=a\,dx^I\)、\(\beta=b\,dx^J\)。
系数乘积求导给出 \(b\,da+a\,db\)，
再把第二项中的 \(db\) 移过 \(p\) 个一次形式，得到因子 \((-1)^p\)。
对所有 \(I,J\) 求和即得结论。
\end{proof}

\subsection{由积分到连续线性泛函}
\label{b1:subsec:D10202}

\begin{definition}\label{b1:def:app-d-current}
从 \(\mathcal D^k(U)\) 到 \(\mathbb R\) 的连续线性泛函称为 \(U\) 上的 \(k\) 维%
\term[liu]{流}[current]，全体记为 \(\mathcal D_k(U)\)。
\symidx{\(\mathcal D_k(U)\)：\(k\) 维流的空间}
若 \(T_j(\omega)\to T(\omega)\) 对每个 \(\omega\in\mathcal D^k(U)\) 成立，
则称 \(T_j\) 弱收敛于 \(T\)，记为 \(T_j\rightharpoonup T\)。
\end{definition}

连续性的含义与分布相同：对每个紧集 \(K\subset U\)，
存在 \(C_K<\infty\) 和非负整数 \(m_K\)，使支集位于 \(K\) 内的形式满足
\[
 |T(\omega)|\leq
 C_K\max_{\substack{|I|=k\\|\alpha|\leq m_K}}
       \|\partial^\alpha\omega_I\|_\infty.
\]
证明可逐系数归结到测试函数空间的连续性。
特别地，零维流就是\chapref{b1:ch:20}的分布，而不只是点质量的线性组合。
这里的“流”也不是依时间演化的映射族；两种对象的英文分别为 current 和 flow。

流的支撑沿用分布的约定：删去所有使该流在其中的测试形式上恒为零的开集，
剩下的相对闭集记为 \(\operatorname{supp}T\)。
若测试形式在此支撑的邻域恒为零，使用光滑划分便知它在 \(T\) 上取值为零。
因此后文可在支撑附近乘截断来解释非紧支集光滑形式的配对。

设 \(M\subset U\) 是有定向的光滑 \(k\) 维子流形，且在紧集内面积有限。
以单位简单切向 \(k\) 向量 \(\tau_M\) 表示定向，定义
\[
 [M](\omega)=\int_M\langle\omega(x),\tau_M(x)\rangle\,d\mathcal H^k(x).
\]
在每个共同紧支集上，积分由系数的一致范数控制，所以这确实是流。
改变取向得到 \(-[M]\)。
从 \(a\) 指向 \(b\) 的线段对应
\[
 [a,b](\omega)
 =\int_0^1\langle\omega(a+t(b-a)),b-a\rangle\,dt.
\]
这个记号表示有向线段流，本附录中不表示无向闭区间。

\begin{definition}\label{b1:def:app-d-current-boundary}
对 \(k\geq1\) 和 \(T\in\mathcal D_k(U)\)，定义其%
\term[liu-bianjie]{边界}[boundary]为
\[
 (\partial T)(\eta)\coloneqq T(d\eta),
 \quad\eta\in\mathcal D^{k-1}(U).
\]
对零维流规定 \(\partial T=0\)。
\symidx{\(\partial T\)：流的边界}
\end{definition}

由外微分的连续性，\(\partial T\) 是 \(k-1\) 维流；
由 \(d^2=0\)，有 \(\partial^2T=0\)。
若 \(T_j\rightharpoonup T\)，则对固定 \(\eta\)，
\(\partial T_j(\eta)=T_j(d\eta)\to T(d\eta)\)，故边界也弱收敛。
边界条件因而能在弱极限中保留，且不需要先知道极限曲面是否光滑。

在线段例子中，一维微积分基本定理给出
\[
 \partial[a,b]=\delta_b-\delta_a.
\]
这同时确定了本书的符号约定：终点为正，起点为负。
对光滑带边界曲面，经典 Stokes 公式写作
\(\partial[M]=[\partial M]\)，边界取“外法向在前”的定向。
流的定义把这个公式推广为边界算子的定义，并不把一切流预先视为某张光滑曲面。

\subsection{余质量与质量范数}
\label{b1:subsec:D10203}

\appref{b1:app:B}在外空间中使用使标准外积基正交归一的欧氏范数，记为 \(|\xi|\)。
当 \(\xi\) 为简单向量时，它等于对应平行多面体的体积；
一般向量可能不是一个平行多面体，需要另一个范数。

\begin{definition}\label{b1:def:app-d-comass}
对交错 \(k\) 次形式 \(\alpha\)，定义其%
\term[yu-zhiliang]{余质量}[comass]为
\[
 \|\alpha\|_{\mathrm{co}}
  \coloneqq\sup\{\,|\langle\alpha,\xi\rangle|:
                   \xi\text{ 为简单 }k\text{ 向量},\ |\xi|=1\,\}.
\]
在 \(k=0\) 时，规定它为绝对值。
外空间上的对偶范数记为
\[
 |\xi|_{\mathrm m}
    \coloneqq\sup_{\|\alpha\|_{\mathrm{co}}\leq1}|\langle\alpha,\xi\rangle|,
\]
称为外向量的质量范数。
\end{definition}

\begin{proposition}\label{b1:prop:app-d-mass-comass}
上述两个量都是范数。对简单向量 \(\xi\)，有 \(|\xi|_{\mathrm m}=|\xi|\)；
对任意 \(\xi\)，则有
\[
 |\xi|\leq|\xi|_{\mathrm m}\leq
           \sqrt{\binom nk}\,|\xi|.
\]
\end{proposition}
\begin{proof}
简单标准基张成整个外空间，因而余质量只在零形式处消失；
三角不等式和齐次性由上确界定义得到。
其对偶量也是有限维空间上的范数。
若 \(\xi\) 简单，取向后可写为
\(\xi=|\xi|\,v_1\wedge\cdots\wedge v_k\)，其中 \(v_i\) 标准正交。
定义本身给出 \(|\xi|_{\mathrm m}\leq|\xi|\)；
取相应的体积形式 \(v_1^*\wedge\cdots\wedge v_k^*\)，
其余质量为一且在该向量上取值 \(|\xi|\)，得到反向不等式。

对一般 \(\xi=\sum_I\xi_Ie_I\)，欧氏单位形式的余质量不超过一，
故欧氏对偶范数不超过质量范数。
另一方面，三角不等式给出
\(|\xi|_{\mathrm m}\leq\sum_I|\xi_I|
\leq\sqrt{\binom nk}\,|\xi|\)。
\end{proof}

这两个范数不能在所有维数上混用。例如在 \(\bigwedge^2\mathbb R^4\) 中，
\[
 \xi=e_1\wedge e_2+e_3\wedge e_4
\]
的欧氏范数为 \(\sqrt2\)，质量范数却为 \(2\)。
形式 \(dx^1\wedge dx^2+dx^3\wedge dx^4\) 的余质量为一，
因为它在正交单位向量 \(v,w\) 上的值可写成 \(\langle Jv,w\rangle\)，其中 \(J\) 为正交变换；
它在 \(\xi\) 上取值二。反向上界来自两项分解。

\begin{definition}\label{b1:def:app-d-current-mass}
流 \(T\in\mathcal D_k(U)\) 的%
\term[liu-zhiliang]{质量}[mass]定义为
\[
 \mathbf M_U(T)\coloneqq
 \sup\{\,|T(\omega)|:\omega\in\mathcal D^k(U),\
                           \sup_x\|\omega(x)\|_{\mathrm{co}}\leq1\,\}.
\]
在 \(\mathbb R^n\) 上简记为 \(\mathbf M(T)\)。
\symidx{\(\mathbf M(T)\)：流的质量}
\end{definition}

对任意开集 \(W\subset U\)，把测试形式限制为支集包含于 \(W\)，便得到局部质量。
质量可以为无穷。
本书采用余质量的对偶定义；Simon 在《Lectures on Geometric Measure Theory》%
\cite[26.6 REMARK]{Simon1983Lectures}明确说明后改用欧氏形式范数。
两种约定在可整流流上给出同一面积，但在一般流上不必数值相等。
有限维范数等价保证局部有界性和相应紧性假设相同，不能由此声称所有质量数值相同。

\begin{proposition}\label{b1:prop:app-d-mass-representation}
若 \(T\) 在每个内部紧集附近质量有限，则存在唯一正 Radon 测度 \(\|T\|\)
及一个 \(\|T\|\)-可测外向量场 \(\vec T\)，使
\[
 T(\omega)=\int_U\langle\omega,\vec T\rangle\,d\|T\|,
 \quad |\vec T|_{\mathrm m}=1\quad\|T\|\text{-几乎处处}.
\]
向量场在 \(\|T\|\)-几乎处处意义下唯一，且对每个开集 \(W\subset U\)，
\(\|T\|(W)=\mathbf M_W(T)\)。
\symidx{\(\lVert T\rVert\)：流的质量测度}
\end{proposition}
\begin{proof}
逐个测试标准系数，局部质量界将每个分量延拓为连续紧支集函数上的局部有界泛函。
由实值 Riesz 表示，每个分量是符号 Radon 测度 \(\mu_I\)。
以 \(\lambda=\sum_I|\mu_I|\) 控制这些分量，写
\(\mu_I=v_I\lambda\)，并记 \(v=\sum_Iv_Ie_I\)。
令 \(q=|v|_{\mathrm m}\)，取 \(\|T\|=q\lambda\) 及 \(\vec T=v/q\)；
在 \(q=0\) 处任意赋值，因为该处不承担 \(\|T\|\)。
范数等价保证这仍是局部有限 Radon 测度。

由对偶不等式，所有允许的测试形式都满足
\(|T(\omega)|\leq\int\|\omega\|_{\mathrm{co}}\,d\|T\|\)，
所以局部质量不超过 \(\|T\|(W)\)。
反向可先取 \(W\) 内紧集 \(K\)，再用余质量单位球中的可数稠密集合，
在每点选取使配对逼近 \(q\) 的形式。
将选择截成有限个可测片，并由 Radon 正则性在各片内部取两两不交的紧集，
使被删去部分的 \(\|T\|\) 质量任意小。
在这些紧集的两两不交邻域上取光滑截断，乘以对应的常值形式后相加。
所成测试形式在每点至多有一个非零项，故余质量不超过一，
其积分任意接近 \(\|T\|(K)\)。
对 \(K\) 取上确界便得到等号。
分量测度及其 Radon--Nikodym 密度的唯一性给出最后的唯一性。
\end{proof}

对 Borel 集 \(A\subset U\)，现在可以定义限制
\[
 (T\llcorner A)(\omega)
       =\int_A\langle\omega,\vec T\rangle\,d\|T\|.
\]
不光滑示性函数的限制只在这种测度表示成立后使用。
一般分布不能直接与任意示性函数相乘。

\begin{definition}\label{b1:def:app-d-normal-current}
若 \(\mathbf M(T)+\mathbf M(\partial T)<\infty\)，则称 \(T\) 为%
\term[zhenggui-liu]{正规流}[normal current]。
只要求上述两项在每个内部紧集附近有限时，称为局部正规流。
\end{definition}

“正规”在此不表示光滑，也不表示与法向量有关。
正规性同时控制对象和边界，而不是只要求对象有限质量。
例如在 \(\mathbb R^2\) 取非零光滑紧支集函数 \(\phi\)，令
\[
 T(\omega_1\,dx^1+\omega_2\,dx^2)
      =\int_{\mathbb R^2}\phi\omega_1\,dx.
\]
其边界为密度 \(-\partial_1\phi\) 的零维流，所以 \(T\) 正规。
它的质量却分布在二维体积上，不可能由可数条曲线的长度积分表示。
下一节的可整流性是独立的几何要求。

% !TeX root = ../../Book1.tex
\section{可整流流、积分流与几何运算}
\label{b1:sec:D103}

正规性控制质量和边界质量，却没有规定质量应集中在什么集合上。
若要描述带重数的曲面，应加入可整流性；
若希望一张曲面在极限中仍由整数层组成，还要加入整数重数。
这几个要求逐层增加，不能把“有面积”“有边界”和“有整数层数”当成同一条件。

\subsection{承载集、重数和取向}
\label{b1:subsec:D10301}

\begin{definition}\label{b1:def:app-d-rectifiable-current}
设 \(1\leq k\leq n\)。若一个流可以写成
\[
 T(\omega)=\int_E\theta(x)
                 \langle\omega(x),\tau(x)\rangle\,d\mathcal H^k(x),
\]
其中 \(E\) 可数 \(k\) 维可整流，\(\theta>0\) 可测，
\(\theta\mathcal H^k\llcorner E\) 局部有限，
且 \(\tau(x)\) 几乎处处是近似切空间的可测单位可分解 \(k\) 向量，
则称 \(T\) 为%
\term[kezhengliu-liu]{可整流流}[rectifiable current]，记作 \(T=[E,\theta,\tau]\)。
若另有 \(\theta(x)\in\mathbb N\) 几乎处处成立，则称其具有整数重数。
\end{definition}

本书在“可整流流”中容许实重数；需整数时明确写出。
Simon 的《Lectures on Geometric Measure Theory》%
\cite[\S27]{Simon1983Lectures}专门讨论 integer multiplicity rectifiable currents，
并说明其与 Federer--Fleming 用语的对应。
查阅其他资料时，不能仅凭 rectifiable 一词判断是否已经包含整数条件。

\ProseParagraphBreak

方向的符号由 \(\tau\) 携带，故取 \(\theta>0\) 不损失一般性。
也可以固定一套可测取向，允许 \(\theta\) 有正负号；两种表示给出相同的流。
由\cref{b1:prop:app-d-mass-representation}以及可分解向量的质量范数，
\[
 \|T\|=\theta\mathcal H^k\llcorner E,\quad
 \mathbf M(T)=\int_E\theta\,d\mathcal H^k.
\]
这里的重数并不是把集合真的复制到别处，而是在同一位置把积分计入多次。

\ProseParagraphBreak

零维的约定单独说明：局部整数重数零维流是
\(\sum_iq_i\delta_{x_i}\)，其中 \(q_i\in\mathbb Z\setminus\{0\}\)，
每个紧集内 \(\sum_{x_i\in K}|q_i|<\infty\)。
因每个非零整数的绝对值至少为一，这等价于每个紧集只含有限个非零点。
全域有限质量时，非零点总数有限。

\begin{definition}\label{b1:def:app-d-integral-current}
若 \(T\) 及 \(\partial T\) 都具有整数重数可整流表示，且
\[
 \mathbf M(T)+\mathbf M(\partial T)<\infty,
\]
则称 \(T\) 为%
\term[jifen-liu]{积分流}[integral current]。
我们把 \(\mathbb R^n\) 上全体 \(k\) 维积分流构成的集合记为 \(\mathbf I_k(\mathbb R^n)\)。
\symidx{\(\mathbf I_k(\mathbb R^n)\)：有限质量积分流}
只在内部紧集上要求相应有限性时，使用“局部积分流”。
\end{definition}

这一定义已经把边界可整流性列入条件，不依赖尚未证明的边界可整流定理。
在更深入的理论中，整数重数可整流流只要边界局部质量有限，
边界便自动整数重数可整流；Simon 的 30.3 THEOREM\cite[30.3 THEOREM]{Simon1983Lectures}给出这一结果。
本附录的紧性和变分问题使用上述明示边界条件的定义。

\begin{proposition}\label{b1:prop:app-d-integral-algebra}
对每个 \(0\leq k\leq n\)，\(\mathbf I_k(\mathbb R^n)\)
对加法、取负和整数倍封闭。
若 \(1\leq k\leq n\) 且 \(T\in\mathbf I_k(\mathbb R^n)\)，则
\(\partial T\in\mathbf I_{k-1}(\mathbb R^n)\)。
\end{proposition}
\begin{proof}
零维情形由有限整数系数点质量的表示直接得到。
以下设 \(k\geq1\)。在两个可整流承载集的交集上，
它们的近似切空间对 \(\mathcal H^k\)-几乎处处的点相同。
这可由可整流集的光滑片分解得到：两张同维光滑片若切空间不同，
其交集在该点附近的 \(k\) 维面积密度为零。
因此在公共部分，两套单位定向相同或相反，重数相加或相减后仍是整数。
删去重数为零的部分，并把负号吸收到定向中，便得到和的整数可整流表示。
三角不等式给出质量有限性，边界之和也有同样性质。
取负、整数倍相同处理。
最后 \(\partial^2T=0\)，所以 \(\partial T\) 的边界是零流，
由积分流的定义可知 \(\partial T\in\mathbf I_{k-1}(\mathbb R^n)\)。
\end{proof}

有限个有向 \(k\) 维单纯形的整数线性组合给出最直接的例子。
称这样的流为整数系数多面体流；允许把单纯形换成有向多面体面，经过有限剖分所得类别不变。
相邻面的公共边界按取向相消，所以一个剖分好的平面区域不会在每条内部剖分线上产生额外边界。
在这种对象上，流的加法就是有向链的加法。

\subsection{有限周长集合给出的积分流}
\label{b1:subsec:D10302}

令 \(E\subset\mathbb R^n\) 可测，定义最高维流
\[
 [E](f\,dx^1\wedge\cdots\wedge dx^n)=\int_E f\,dx.
\]
它的质量等于 \(|E|\)。
最高维切空间就是整个 \(\mathbb R^n\)，所以这个表示自然具有重数一和标准取向。

\begin{proposition}\label{b1:prop:app-d-perimeter-current}
若 \(E\) 具有有限体积和有限周长，则 \([E]\) 为积分流，且
\[
 \mathbf M([E])=|E|,\quad
 \|\partial[E]\|=\mathcal H^{n-1}\llcorner\partial^*E.
\]
边界取向 \(\tau_E\) 由
\[
 \nu_E\wedge\tau_E=e_1\wedge\cdots\wedge e_n
\]
确定，其中 \(\nu_E\) 为集合的测度论外法向。
\end{proposition}
\begin{proof}
每个 \((n-1)\) 次形式都唯一写成
\[
 \eta=\sum_{i=1}^n(-1)^{i-1}X_i\,
        dx^1\wedge\cdots\wedge\widehat{dx^i}\wedge\cdots\wedge dx^n
\]
其中 \(X\) 为光滑向量场，帽号表示删去该项。
直接外微分得到 \(d\eta=(\operatorname{div}X)\,dx^1\wedge\cdots\wedge dx^n\)。
由\cref{b1:thm:gauss-green-finite-perimeter}，
\[
 \partial[E](\eta)
   =\int_E\operatorname{div}X\,dx
   =\int_{\partial^*E}X\cdot\nu_E\,d\mathcal H^{n-1}.
\]
上述取向满足
\(\langle\eta,\tau_E\rangle=X\cdot\nu_E\)，故
\(\partial[E]=[\partial^*E,1,\tau_E]\)。
由 De Giorgi 结构定理，\(\partial^*E\) 可数 \((n-1)\) 维可整流，
其面积等于周长；两个质量均有限。
\end{proof}

从这个例子可看出，有界变差函数理论已经包含一种重要的流。
集合的拓扑边界可能因修改零测集而完全改变，流 \([E]\) 及其边界却不会改变。
边界也不应写成整个拓扑边界上的面积积分，正确承载集是约化边界。

\subsection{光滑映射下的推前}
\label{b1:subsec:D10303}

设 \(f:U\to V\) 光滑。
外代数的拉回给出
\[
 (f^*\omega)_x(\xi)
   =\omega_{f(x)}\bigl((\bigwedge\nolimits^kDf_x)\xi\bigr).
\]
若每个目标紧集的逆像与 \(\operatorname{supp}T\) 相交后仍紧，则定义
\[
 (f_\#T)(\omega)=T(f^*\omega).
\]
当拉回不在全域紧支集时，右端在它与 \(\operatorname{supp}T\) 相交的紧集附近乘一个等于一的截断。
不同截断的差在支撑附近为零，因而给出同一数值。
本附录后面主要使用紧支集流，此时连续映射自动满足所需的支撑条件。

\begin{proposition}\label{b1:prop:app-d-pushforward}
在上述支撑条件下，推前是流，且
\[
 \partial(f_\#T)=f_\#(\partial T).
\]
若 \(T\) 质量有限且 \(Df\) 在其支撑上有界，则
\[
 \mathbf M(f_\#T)\leq
   \left(\sup_{\operatorname{supp}T}\|Df\|_{\mathrm{op}}\right)^k
          \mathbf M(T).
\]
若 \(T\) 为整数重数可整流流，则 \(f_\#T\) 也具有整数重数可整流表示；
若 \(T\) 为紧支集积分流，则 \(f_\#T\) 为积分流。
\end{proposition}
\begin{proof}
链式法则保证拉回在固定紧支集上的各阶导数受控，所以推前连续。
逐坐标计算给出 \(d(f^*\eta)=f^*(d\eta)\)：
一阶系数满足通常链式法则，二阶偏导项因反对称性相消。
于是
\[
 \partial(f_\#T)(\eta)=T(f^*d\eta)
       =T(d f^*\eta)=f_\#(\partial T)(\eta).
\]
对单位可分解向量 \(\xi\)，其像仍可分解，且
\(|(\bigwedge^kDf)\xi|\leq\|Df\|_{\mathrm{op}}^k\)。
由此
\(\|f^*\omega\|_{\mathrm{co}}\leq
\|Df\|_{\mathrm{op}}^k\cdot\|\omega\circ f\|_{\mathrm{co}}\)，给出质量估计。

若 \(T=[E,\theta,\tau]\)，则
\[
 f_\#T(\omega)
   =\int_E\theta(x)
       \left\langle\omega(f(x)),(\bigwedge\nolimits^kd^Ef_x)\tau(x)\right\rangle
          \,d\mathcal H^k(x).
\]
对切向 Jacobi 因子正的部分应用可整流集上的面积公式。
在几乎每个像点，所有非退化原像切空间的像等于该像集的近似切空间；
各项取向因此只差一个正负号。
选定像集的一套可测取向后，其重数为各原像整数重数的带符号和。
质量估计保证这些绝对值的和几乎处处有限，所以该和仍为整数。
切向 Jacobi 因子为零的部分不给积分贡献。
对 \(T\) 及 \(\partial T\) 分别使用该结论，再结合边界交换式，即得最后一句。
\end{proof}

像的几何面积不能总用带符号重数的绝对值之前的求和来代替。
两片反向曲面映到同一位置时会相消，故质量估计一般只是“不超过”。
Lipschitz 映射也有推前理论，但一般流的测试形式不能直接拉回成光滑形式。
本节先使用光滑版本；在建立积分流紧性后，%
\cref{b1:prop:app-d-lipschitz-pushforward}将利用同伦控制完成
紧支集积分流上的 Lipschitz 延拓。

\subsection{乘积与同伦公式}
\label{b1:subsec:D10304}

设 \(T\) 为紧支集正规 \(k\) 维流。
在 \(\mathbb R\times\mathbb R^n\) 中，把测试 \((k+1)\) 次形式写成
\(\Omega=dt\wedge\alpha(t,x)+\beta(t,x)\)，
其中 \(\alpha\) 不含 \(dt\)，\(\beta\) 也不含 \(dt\)。
定义
\[
 ([0,1]\times T)(\Omega)=\int_0^1T\bigl(\alpha(t,\cdot)\bigr)\,dt.
\]
该积分的绝对值受质量界控制，且支撑包含在
\([0,1]\times\operatorname{supp}T\) 内，因此定义了有限质量流。

\begin{proposition}\label{b1:prop:app-d-homotopy}
设 \(h:[0,1]\times\mathbb R^n\to\mathbb R^N\)
在该闭条带的邻域上光滑，并记 \(h_t(x)=h(t,x)\)。
对紧支集正规流 \(T\)，有
\[
 \partial h_\#([0,1]\times T)
  =(h_1)_\#T-(h_0)_\#T-h_\#([0,1]\times\partial T).
\]
\end{proposition}
\begin{proof}
先验证乘积边界。
把一个 \(k\) 次测试形式写成
\(\eta=dt\wedge\gamma(t,x)+\alpha(t,x)\)。
其外微分中含 \(dt\) 的部分为
\[
 dt\wedge(\partial_t\alpha-d_x\gamma).
\]
所以
\[
 \begin{split}
 \partial([0,1]\times T)(\eta)
 &=\int_0^1 T(\partial_t\alpha)\,dt
     -\int_0^1\partial T(\gamma)\,dt\\
 &=T(\alpha(1,\cdot))-T(\alpha(0,\cdot))
     -([0,1]\times\partial T)(\eta).
 \end{split}
\]
第一项积分使用一维微积分基本定理；通过质量表示可直接把 \(T\) 与积分交换。
然后应用推前与边界交换的结论，即得公式。
\end{proof}

这条公式说明，两个推前之间的差不一定本身就是一条边界：
当原流有边界时，边界移动扫出的侧面还要计入。
在估计平坦范数时遗漏这项，就会错误地只用 \(\mathbf M(T)\) 控制所有平移。

\ProseParagraphBreak

特别地，令 \(h(t,x)=a+t(x-a)\)。
若 \(T\) 是紧支集积分 \(k\) 维流且 \(k\geq1\)，则
\((h_0)_\#T=0\)；
若又有 \(\partial T=0\)，便得到一个填充
\[
 C_aT=h_\#([0,1]\times T),\quad \partial C_aT=T.
\]
乘积具有整数重数可整流表示，光滑推前保持该性质，因此 \(C_aT\) 为积分流。
若 \(\operatorname{supp}T\subset\overline{B(a,R)}\)，
其切向外积在时间 \(t\) 的长度至多为 \(R t^k\)，从而
\[
 \mathbf M(C_aT)\leq\frac{R}{k+1}\,\mathbf M(T).
\]
在零维时，常值推前等于总系数乘 \(\delta_a\)，所以只有总系数为零的零维流才能这样填充。
例如 \(\delta_b-\delta_c\) 有线段填充，而单个 \(\delta_b\) 没有紧支集的一维填充。

% !TeX root = ../../Book1.tex
\section{切片与边界的局部化}
\label{b1:sec:D104}

把一张有边界的曲面沿一个水平面切开，会出现两种边界：
原边界留在下方的部分，以及切割新产生的截面。
后者并不是把原流直接限制在水平集上。
一张二维曲面的面积测度通常不给一条水平截线任何质量，
但截线本身可以有正长度。切片必须降低维数，不能只乘一个示性函数。

\subsection{由边界差定义切片}
\label{b1:subsec:D10401}

本节先设 \(T\) 是 \(\mathbb R^n\) 上的紧支集正规 \(k\) 维流，\(k\geq1\)，
且 \(f:\mathbb R^n\rightarrow\mathbb R\) 为 Lipschitz 函数。
有限质量表示已经允许对 Borel 集作限制，所以
\[
 S_t=\partial(T\llcorner\{f<t\})
             -(\partial T)\llcorner\{f<t\}
\]
对每个实数 \(t\) 都定义了一个流。
此时尚不知道其质量是否有限。
对几乎所有水平值成立的有限质量结论，是切片理论首先需要证明的内容。

\begin{theorem}\label{b1:thm:app-d-normal-slicing}
在上述条件下，\(S_t\) 对几乎每个 \(t\) 是正规 \(k-1\) 维流，且
\[
 \int_{\mathbb R}\mathbf M(S_t)\,dt
       \leq\operatorname{Lip}(f)\,\mathbf M(T).
\]
其支撑包含于 \(\operatorname{supp}T\cap\{f=t\}\)，并满足
\[
 \partial S_t=-\langle\partial T,f,t\rangle
 \quad\text{对几乎每个 }t.
\]
右端为对 \(\partial T\) 使用同一构造的切片；零维流的切片约定为零。
\end{theorem}

在这些良好水平上，将 \(S_t\) 记为 \(\langle T,f,t\rangle\)，称为 \(T\)
由 \(f\) 在水平 \(t\) 处确定的%
\term[liu-qiepian]{切片}[slice]。
\symidx{\(\langle T,f,t\rangle\)：流的标量切片}
同一术语用于多分量映射时，还须逐个固定坐标顺序；本节只需标量版本。

\begin{proof}
先假设 \(f\) 光滑且 \(|Df|\leq L\)。
取光滑递增函数 \(\chi\)，使 \(0\leq\chi\leq1\)，
在 \((-\infty,0]\) 上为零、在 \([1,\infty)\) 上为一。
令
\[
 g_{\varepsilon,t}(x)
       =\chi\!\left(\frac{t-f(x)}{\varepsilon}\right).
\]
乘以非紧支集函数时，在 \(T\) 的固定紧支撑附近使用截断。
由外微分的乘积法则，对测试 \(k-1\) 次形式 \(\eta\)，有
\[
 \begin{split}
 S_{\varepsilon,t}(\eta)
 &=\bigl[\partial(T\llcorner g_{\varepsilon,t})
               -(\partial T)\llcorner g_{\varepsilon,t}\bigr](\eta)\\
 &=-T(dg_{\varepsilon,t}\wedge\eta)\\
 &=T\!\left(\frac1\varepsilon
       \chi'\!\left(\frac{t-f}{\varepsilon}\right)df\wedge\eta\right).
 \end{split}
\]
对一次形式 \(a\) 和交错形式 \(\eta\)，有
\(\|a\wedge\eta\|_{\mathrm{co}}\leq |a|\cdot\|\eta\|_{\mathrm{co}}\)。
验证时可把任一单位简单 \(k\) 向量的第一个正交基方向选成
\(a\) 在其张成空间上的投影方向；其余方向被 \(a\) 消去，
只剩一个系数不超过 \(|a|\) 的配对。
于是
\[
 \mathbf M(S_{\varepsilon,t})
 \leq L\int\frac1\varepsilon
           \chi'\!\left(\frac{t-f(x)}{\varepsilon}\right)\,d\|T\|(x).
\]
对 \(t\) 积分，Tonelli 定理及 \(\int\chi'=1\) 给出
\(\int\mathbf M(S_{\varepsilon,t})\,dt\leq L\mathbf M(T)\)。

测度 \(f_\#(\|T\|+\|\partial T\|)\) 有限，故它的原子至多可数。
排除对应水平后，\(g_{\varepsilon,t}\) 对这两个测度几乎处处趋于
\(\mathbf1_{\{f<t\}}\)。
控制收敛说明两个被限制的流分别按质量收敛，
从而 \(S_{\varepsilon,t}\rightharpoonup S_t\)。
质量是对测试形式所得连续函数的上确界，所以对弱收敛下半连续。
因此 Fatou 引理给出
\[
 \int\mathbf M(S_t)\,dt\leq L\mathbf M(T).
\]
这里的被积函数可测：在固定支撑邻域内取可数个光滑形式，
使它们在函数及一阶导数的一致范数下稠密，
便可把质量写成可数上确界。对每个这样的形式，限制积分关于 \(t\) 可测。

对一般 Lipschitz 函数，取全空间磨光 \(f_j\)，使 \(f_j\rightarrow f\) 一致且
\(\operatorname{Lip}(f_j)\leq\operatorname{Lip}(f)\)。
在上面排除的水平以外，
\(\mathbf1_{\{f_j<t\}}\rightarrow\mathbf1_{\{f<t\}}\)
对 \(\|T\|+\|\partial T\|\) 几乎处处成立。
相应的边界差因而弱收敛于 \(S_t\)。
再用质量下半连续性与 Fatou 引理，得到同一个积分上界。
这个论证只拉回了函数的水平集，没有把不光滑形式当作一般分布的测试形式。

若一个测试形式的支撑位于 \(\{f<t\}\)，两个限制在该支撑附近都等于原流，
故边界差为零；位于 \(\{f>t\}\) 时两者都为零。
对 \(\{f\neq t\}\) 内的测试形式使用有限光滑划分，得到所述支撑包含。
最后对定义式再取边界，利用 \(\partial^2T=0\)，有
\[
 \partial S_t=-\partial\bigl((\partial T)\llcorner\{f<t\}\bigr)
             =-\langle\partial T,f,t\rangle.
\]
若 \(k\geq2\)，对正规流 \(\partial T\) 使用已证质量估计，其切片几乎处处有限质量；
若 \(k=1\)，零维流的边界为零。两种情形都得到 \(S_t\) 的正规性。
\end{proof}

定理只保证几乎处处良好，不能要求每个预先指定的水平都具有同样性质。
不过，若需在一个小区间里选择合适水平，积分上界总能提供它。
例如对长度为 \(\delta\) 的区间 \(J\)，可以选取 \(t\in J\)，使
\[
 \mathbf M(\langle T,f,t\rangle)
       \leq \frac{\operatorname{Lip}(f)}{\delta}\,\mathbf M(T).
\]
这一选择常用于先截取紧支撑部分，再控制新产生的边界。

\subsection{可整流流的几何截面}
\label{b1:subsec:D10402}

设 \(T=[E,\theta,\tau]\)，且 \(f\) 光滑。
在 \(E\) 的近似切空间上，把 \(Df\) 的切向部分按\chapref{b1:ch:19}记为 \(d^Ef\)，
其对应切向梯度记为 \(\nabla_Ef\)。
当该梯度非零时，截面切空间是
\[
 \operatorname{Tan}^k(E,x)\cap\ker Df_x.
\]
若 \(v=\nabla_Ef/|\nabla_Ef|\)，则截面取向 \(\tau_t\) 由
\(v\wedge\tau_t=\tau\) 确定。这是“朝 \(f\) 增大方向在前”的约定。

\begin{theorem}\label{b1:thm:app-d-rectifiable-slicing}
若 \(T=[E,\theta,\tau]\) 为紧支集积分 \(k\) 维流，且 \(f\) 光滑、梯度有界，
则几乎每个 \(\langle T,f,t\rangle\) 是积分流。
它由
\[
 E_t=E\cap\{f=t\}\cap\{|\nabla_Ef|>0\}
\]
承载，具有重数 \(\theta|_{E_t}\) 和上述取向 \(\tau_t\)。
并且
\[
 \int_{\mathbb R}\mathbf M(\langle T,f,t\rangle)\,dt
        =\int_E\theta\,|\nabla_Ef|\,d\mathcal H^k.
\]
\end{theorem}
\begin{proof}
可整流集上的余面积公式给出：几乎每个 \(E_t\) 可数 \(k-1\) 维可整流，
并具有上述切空间。
这可在可数光滑片上应用隐函数定理，再删去切向微分退化部分得到；
退化部分的截面面积积分由余面积公式等于零。
把这些片合并，所得截面流
\[
 R_t(\eta)=\int_{E_t}\theta\langle\eta,\tau_t\rangle\,d\mathcal H^{k-1}
\]
在几乎每个水平具有有限质量，且重数仍为整数。

前一证明中的 \(S_{\varepsilon,t}\) 可写成
\[
 S_{\varepsilon,t}(\eta)
    =\int_{\mathbb R}\frac1\varepsilon
          \chi'\!\left(\frac{t-s}{\varepsilon}\right)R_s(\eta)\,ds.
\]
这里使用配对恒等式
\(\langle df\wedge\eta,\tau\rangle
 =|\nabla_Ef|\langle\eta,\tau_t\rangle\)，
以及余面积公式。
对每个固定测试形式，右边是可积函数 \(s\mapsto R_s(\eta)\) 的单侧近似恒等算子。
Lebesgue 微分定理说明它几乎处处趋于 \(R_t(\eta)\)。
取一个可数、在一阶一致范数下稠密的测试族，
把这些例外集与上一证明的例外集合并，
便在同一满测水平集上得到 \(S_t=R_t\)。
从稠密测试族推广时，\(S_t\) 和 \(R_t\) 已有有限质量，不需要再为每个形式选择子列。

当 \(k\geq2\) 时，对 \(\partial T\) 使用同一论证，并利用切片边界公式，
可知 \(\partial S_t\) 也具有整数重数可整流表示。
当 \(k=1\) 时，\(S_t\) 为零维整数重数流，其边界为零。
所以 \(S_t\) 为积分流。
最后，简单单位截面取向的质量范数为一，
对 \(\theta\) 应用余面积公式，即得等式。
\end{proof}

Simon 的《Lectures on Geometric Measure Theory》%
\cite[28.1 LEMMA, 28.2 REMARK, 28.3, 28.4 DEFINITION, and 28.5 LEMMA]{Simon1983Lectures}给出 Lipschitz 函数的同一几何描述。
其推广还要在可整流集的光滑片上处理切向微分及函数逼近；
不能把一般流的形式拉回直接代入来替代这一步。
上面的完整几何证明采用光滑函数；正规流的边界差构造与积分估计已经覆盖 Lipschitz 函数。
后文按网格坐标切割只需光滑的线性坐标函数。

\subsection{符号、重数与切割产生的边界}
\label{b1:subsec:D10403}

取带标准定向的矩形
\(Q=(0,a)\times(0,b)\subset\mathbb R^2\)，令 \(T=[Q]\)、\(f(x,y)=x\)。
当 \(0<t<a\) 时，切片是从 \((t,0)\) 指向 \((t,b)\) 的线段。
其取向为 \(e_2\)，因为 \(e_1\wedge e_2\) 是矩形的取向。
切片的两个端点满足
\[
 \partial\langle T,f,t\rangle
    =\delta_{(t,b)}-\delta_{(t,0)}.
\]
原矩形上边向左、下边向右，因此
\[
 \langle\partial T,f,t\rangle
    =\delta_{(t,0)}-\delta_{(t,b)}.
\]
两式直接展示了边界交换式中的负号。若把原流乘以整数 \(q\)，截面也乘以 \(q\)，
维数降低不会消除重数。

另一方面，\(\|T\|(\{x=t\})=0\)，故作为二维流的限制
\(T\llcorner\{x=t\}\) 等于零。
这是切片不等于水平集限制的具体例子。
对高度函数 \(f(x,y)=x+y\)，截面长度随 \(t\) 改变，
但其总量满足
\[
 \int_{\mathbb R}\mathbf M(\langle T,f,t\rangle)\,dt
       =\sqrt2\,ab,
\]
因为 \(|\nabla f|=\sqrt2\)。额外的 \(\sqrt2\) 来自水平参数的速度，
不是取向或重数造成的面积修正。

若 \(T\) 原本没有边界，切割后的新边界就是
\(\partial(T\llcorner\{f<t\})=\langle T,f,t\rangle\)。
若原流有边界，两部分必须同时保留。
这一恒等式把可控切割接到下一节的网格变形上：
要把曲面移到低维骨架，既要估计移动扫出的体积，也要估计边界在切割和移动中产生的误差。

% !TeX root = ../../Book1.tex
\section{平坦范数、变形与积分流紧性}
\label{b1:sec:D105}

两张靠得很近但互不重合的曲面，作为向量测度可以相差很大的总变差。
若它们之间只夹着一层很薄的区域，几何上却应当接近。
合适的距离允许把差拆成两部分：一部分直接按质量计价，另一部分作为某个高一维流的边界，
按这个填充的质量计价。这样既能测量小块误差，也能测量曲面移动所扫过的薄层。

\subsection{允许填充的距离}
\label{b1:subsec:D10501}

\begin{definition}\label{b1:def:app-d-flat-norm}
对 \(\mathbb R^n\) 上的 \(k\) 维流 \(T\)，定义%
\term[pingtan-fanshu]{平坦范数}[flat norm]为
\[
 \mathbb F(T)=\inf\{\,\mathbf M(A)+\mathbf M(B):
       T=A+\partial B,\ \mathbf M(A),\mathbf M(B)<\infty\,\},
\]
其中 \(A\) 为 \(k\) 维流，\(B\) 为 \(k+1\) 维流。
若不存在这种分解，则取值为 \(+\infty\)。
在最高维时 \(B=0\)。这里允许实系数的有限质量流，不要求 \(A,B\) 是积分流。
\symidx{\(\mathbb F(T)\)：实系数平坦范数}
\end{definition}

定义里没有要求下确界能取到，也没有限定填充的支撑。
后面讨论固定紧集内的对象时，紧性假设作用于原流，而不是静默修改本定义的分解集合。

\begin{proposition}\label{b1:prop:app-d-flat-properties}
在平坦范数有限的流组成的向量空间上，\(\mathbb F\) 是范数，且
\[
 \mathbb F(T)\leq\mathbf M(T),\quad
 \mathbb F(\partial T)\leq\mathbb F(T).
\]
对任意测试 \(k\) 次形式 \(\omega\)，有
\[
 |T(\omega)|
 \leq
 \max\!\left\{\sup_x\|\omega(x)\|_{\mathrm{co}},
               \sup_x\|d\omega(x)\|_{\mathrm{co}}\right\}\mathbb F(T).
\]
因而平坦范数收敛蕴含流的弱收敛。
\end{proposition}
\begin{proof}
分解相加、乘以标量，分别给出三角不等式和绝对齐次性；
齐次性的反向由除以非零标量得到。
取 \(A=T,B=0\)，得到第一个上界。
若 \(T=A+\partial B\)，则 \(\partial T=\partial A\)，
故 \(\mathbb F(\partial T)\leq\mathbf M(A)\)。
对所有分解取下确界，得到第二个上界。

对同一个分解，
\[
 |T(\omega)|=|A(\omega)+B(d\omega)|
 \leq \mathbf M(A)\sup_x\|\omega(x)\|_{\mathrm{co}}
       +\mathbf M(B)\sup_x\|d\omega(x)\|_{\mathrm{co}}.
\]
取下确界即得测试估计。
若 \(\mathbb F(T)=0\)，则所有测试形式上的值都为零，所以 \(T=0\)。
这也证明平坦范数收敛必然弱收敛。
\end{proof}

Simon 的《Lectures on Geometric Measure Theory》%
\cite[\S31]{Simon1983Lectures}在整数重数流之间采用要求分解项也为整数重数的局部距离。
其允许的分解较少，不能不加说明地把定义中的下确界照搬为本书的 \(\mathbb F\)。
下文会直接证明实系数平坦范数在正规流有界族上的紧性，
再单独使用整数闭合定理保留积分流结构。

例如，平行的两个单位正方形取相同取向，间距为 \(h>0\)。
它们之差的质量是 \(2\)，与 \(h\) 无关；
但借助夹层和沿边界扫出的侧面，平坦范数至多为常数乘 \(h\)。
有边界的曲面移动时，夹层的边界不仅含上下两面，侧面项不能漏去。
下面用同伦公式给出统一估计。

\begin{proposition}\label{b1:prop:app-d-translation-flat}
设 \(T\) 为紧支集正规流，\(\tau_h(x)=x+h\)。则
\[
 \mathbb F\bigl((\tau_h)_\#T-T\bigr)
       \leq |h|\bigl(\mathbf M(T)+\mathbf M(\partial T)\bigr).
\]
\end{proposition}
\begin{proof}
设 \(H_h(t,x)=x+th\)，并记
\(A_hT=(H_h)_\#([0,1]\times T)\)。
由同伦公式，
\[
 (\tau_h)_\#T-T=\partial(A_hT)+A_h(\partial T).
\]
沿时间方向的导数是 \(h\)，其他方向的导数是恒等映射。
由余质量对偶估计，
\[
 \mathbf M(A_hT)\leq |h|\,\mathbf M(T),\quad
 \mathbf M(A_h\partial T)\leq |h|\,\mathbf M(\partial T).
\]
把两项作为平坦范数定义中的分解，即得结论。
\end{proof}

\subsection{正规流的解析紧性}
\label{b1:subsec:D10502}

质量控制首先带来一个向量测度的弱极限。这里的困难尚不在整数，
而在选择同一条子列，使所有测试形式同时收敛。

\begin{proposition}\label{b1:prop:app-d-normal-weak-compactness}
设 \(K\subset\mathbb R^n\) 紧，\(\operatorname{supp}T_j\subset K\)，且
\[
 \sup_j\bigl(\mathbf M(T_j)+\mathbf M(\partial T_j)\bigr)<\infty.
\]
则有子列弱收敛于支撑包含于 \(K\) 的正规流 \(T\)，并且
\[
 \mathbf M(T)\leq\liminf_j\mathbf M(T_j),\quad
 \mathbf M(\partial T)\leq\liminf_j\mathbf M(\partial T_j).
\]
\end{proposition}
\begin{proof}
将 \(T_j\) 的有限质量表示按标准外积基分解。
每个分量都是支撑在 \(K\) 上、总变差一致有界的符号测度。
空间 \(C(K)\) 可分；对其中一个可数稠密族逐次抽取并取对角子列，
使所有分量在该族上同时收敛。
一致总变差界把收敛推广到整个 \(C(K)\)，
Riesz 表示给出极限分量测度。
将它们组合，得到支撑在 \(K\) 上的有限质量流 \(T\)，且 \(T_j\rightharpoonup T\)。

对任意测试 \(k-1\) 次形式 \(\eta\)，
\[
 \partial T_j(\eta)=T_j(d\eta)\longrightarrow T(d\eta)=\partial T(\eta).
\]
原有边界质量界故也约束这个极限，说明 \(T\) 正规。
最后，固定任一个余质量不超过一的测试形式，先取极限再用质量上界，
对所有测试形式取上确界，得到两条下半连续不等式。
对远离 \(K\) 的测试形式，原序列恒为零，故极限仍以 \(K\) 为支撑范围。
\end{proof}

接下来证明，在这样的有界族里，弱收敛还能升级为本节的平坦范数收敛。
与\chapref{b1:ch:25}的平移紧性准则相似，证明先把对象统一磨光，
再把小尺度误差与固定尺度下的紧性分开。

\begin{proposition}\label{b1:prop:app-d-normal-weak-flat}
在前一命题的假设下，若整列 \(T_j\rightharpoonup T\)，则
\(\mathbb F(T_j-T)\rightarrow0\)。
\end{proposition}
\begin{proof}
取非负光滑磨光核 \(\rho_\varepsilon\)，支撑于 \(B(0,\varepsilon)\)，单位积分。
对有限质量流逐系数卷积，等价地定义
\[
 T_\varepsilon
    =\int_{\mathbb R^n}\rho_\varepsilon(h)(\tau_h)_\#T\,dh.
\]
这个积分可按每个测试形式理解，也可按有限维向量测度逐分量理解。
将上一命题证明中的 \(A_hT\) 逐分量积分，记所得流为 \(B_\varepsilon T\)。
其质量至多为 \(\varepsilon\mathbf M(T)\)，支撑包含于 \(K+\overline{B(0,\varepsilon)}\)，
并且边界与积分交换。因此
\[
 T_\varepsilon-T
       =\partial(B_\varepsilon T)+B_\varepsilon(\partial T),\quad
 \mathbb F(T_\varepsilon-T)
       \leq\varepsilon\bigl(\mathbf M(T)+\mathbf M(\partial T)\bigr).
\]
同样的估计对全部 \(T_j\) 使用同一个常数成立。

固定 \(\varepsilon>0\)。
流 \((T_j)_\varepsilon\) 的各分量现在具有光滑密度，
其函数值和一阶导数由原序列的质量界及 \(\rho_\varepsilon\) 的相应导数一致控制。
弱收敛使每个固定点的密度值收敛到 \(T_\varepsilon\) 的密度：
卷积核关于原变量是光滑测试函数，在 \(K\) 附近乘一个固定截断即可。
共同紧支撑及等度连续性把逐点收敛提升为一致收敛。
具体地，在紧支撑上取有限的充分细网格，先控制各网格点的差，再用共同连续模控制网格间隙。
故
\[
 \mathbf M\bigl((T_j)_\varepsilon-T_\varepsilon\bigr)\longrightarrow0.
\]
这里可以先估计每个系数的 \(L^1\) 范数，再用有限维质量范数等价得到上述结论。

三角不等式给出
\[
 \limsup_j\mathbb F(T_j-T)\leq2C\varepsilon
\]
其中 \(C\) 为原序列与极限的共同正规质量界。
最后令 \(\varepsilon\downarrow0\)，得到所需收敛。
\end{proof}

这两个解析命题对实系数正规流成立，但没有证明极限是一张有整数层数的曲面。
例如前面给出的光滑体积密度一维流本来就是正规流。
要把这一紧性用于曲面的面积极小化，必须补上整数重数可整流类的闭合性。

\subsection{网格变形与离散逼近}
\label{b1:subsec:D10503}

整数条件在固定网格上特别清楚。
给定边长 \(\delta\) 的立方网格，一个整数系数 \(k\) 维多面体流
\[
 P=\sum_F q_F[F],\quad q_F\in\mathbb Z
\]
若不为零，其质量至少为 \(\delta^k\)。
这里对网格的 \(k\) 维面 \(F\) 取固定取向，求和只有有限个非零项，
不同面的相对内部互不相交。这个正的最小质量是实系数所不具备的。

将一般流移到网格上并不是作一个处处光滑的最近点投影。
网格低维骨架不存在这种全局投影；正确构造要在各高维胞腔中选择中心，
用径向映射逐层推出，并把移动产生的流和边界误差一起记录。

\begin{theorem}[Federer--Fleming 变形]\label{b1:thm:app-d-deformation}
设 \(1\leq k<n\)，\(T\in\mathbf I_k(\mathbb R^n)\) 有紧支集，\(\delta>0\)。
则存在整数系数网格多面体 \(k\) 维流 \(P\)，以及紧支集积分流
\(R\in\mathbf I_{k+1}(\mathbb R^n)\)、\(S\in\mathbf I_k(\mathbb R^n)\)，使
\[
 T-P=\partial R+S,
\]
并且
\[
 \begin{aligned}
 \mathbf M(P)&\leq C\mathbf M(T),&
 \mathbf M(\partial P)&\leq C\mathbf M(\partial T),\\
 \mathbf M(R)&\leq C\delta\,\mathbf M(T),&
 \mathbf M(S)&\leq C\delta\,\mathbf M(\partial T).
 \end{aligned}
\]
常数 \(C\) 只依赖 \(n,k\)。
流 \(P,R\) 的支撑位于 \(\operatorname{supp}T\) 的 \(C\delta\) 邻域，
而 \(\partial P,S\) 的支撑位于 \(\operatorname{supp}\partial T\) 的 \(C\delta\) 邻域。
\end{theorem}

这条%
\term[federer-fleming-bianxing]{Federer--Fleming 变形定理}[Federer--Fleming deformation theorem]%
属于流理论的主要技术结果。
此处采用 Simon 的《Lectures on Geometric Measure Theory》%
\cite[29.1 THEOREM and 29.3 COROLLARY]{Simon1983Lectures}；该版本对 \(\mathbf M(P)\)
的估计不需要另加 \(\delta\mathbf M(\partial T)\)。
历史原文为 Federer 与 Fleming 的《Normal and Integral Currents》%
\cite{FedererFleming1960Normal}。

\begin{proofsketch}
从 \(n\) 维网格胞腔开始，在每个胞腔内部选一点，把流沿射线推向胞腔边界。
在一个 \(m\) 维胞腔内，径向映射的导数由
\(C\delta/\operatorname{dist}(x,a)\) 控制。
只要 \(m>k\)，对中心 \(a\) 积分时，核
\(\operatorname{dist}(x,a)^{-k}\) 在 \(m\) 维内可积。
Fubini 因而允许选择使 \(k\) 维质量增量有界的中心，
而不是在奇点处直接假设投影的 Lipschitz 常数有限。
对边界也作相容的选择，并逐层推进到 \(k\) 维骨架。

每次变形的同伦扫出一个高一维流；其厚度不超过常数乘 \(\delta\)，
所以总质量由 \(C\delta\mathbf M(T)\) 控制。
原边界的同伦产生同维误差，估计为 \(C\delta\mathbf M(\partial T)\)。
在最终网格面内部，零边界条件使重数为常数；
推前和切片保留整数重数，故最终对象是整数系数网格流。
胞腔内的局部构造还给出所述支撑控制。

这里未展开奇点邻域的切除与极限、各胞腔映射的一致拼接、
以及在最后一层修正边界的步骤。
这些步骤共同承担质量估计和整数性，完整证明见上述第29节；
不能只凭画出一个径向投影就认为变形定理已经证明。
\end{proofsketch}

由分解直接得到
\[
 \mathbb F(T-P)\leq
       C\delta\bigl(\mathbf M(T)+\mathbf M(\partial T)\bigr).
\]
因此积分流可以按平坦范数逼近为整数网格流。
但这并不是质量范数逼近：把一条斜线移到坐标网格上，承载集几乎处处不同，
质量差可以一直不小。网格的作用是控制几何移动后的误差，而不是逐点保留原切向。

\subsection{整数闭合与紧性定理}
\label{b1:subsec:D10504}

网格上的整数系数虽然离散，网格尺度却在趋零。
仅说“整数列的极限仍是整数”不足以处理承载集、切空间和质量都在变化的流。
必须先证明极限重新具有可整流表示，随后才能讨论其中的重数。

\begin{theorem}[Federer--Fleming 整数闭合]\label{b1:thm:app-d-integral-closure}
设 \(U\subset\mathbb R^n\) 开，\(T_j\) 是 \(U\) 内的局部积分 \(k\) 维流，
\(T_j\rightharpoonup T\)，并且对每个 \(W\Subset U\)，有
\[
 \sup_j\bigl(\mathbf M_W(T_j)+\mathbf M_W(\partial T_j)\bigr)<\infty.
\]
则 \(T\) 为局部积分流。
\end{theorem}
\begin{proofsketch}
可按维数归纳。零维时，每个紧集内只有一致有界个带非零整数系数的点，
取点的位置子列并合并重合极限，即保留整数性。
高维时先用切片选取良好球面，使截取流的边界质量受控；
归纳假设使这些低维边界在极限中仍具有整数结构。
用等周填充与局部比较排除极限质量在正密度缺失部分的弥散，
再利用流的可整流性定理得到极限的切向表示。
最后在可整流片上切片或投影，把重数的识别降到整数计数。

这一证明的难点包括从局部填充估计推出正密度，
以及从正密度和边界控制推出切向可整流表示。
这些并不是弱-*紧性本身的后果。Simon 的
《Lectures on Geometric Measure Theory》%
\cite[32.1 THEOREM, 32.2 THEOREM, and 32.3 REMARKS]{Simon1983Lectures}给出完整证明，
其中需要前面的变形、等周和可整流结构理论。
本附录引用这一闭合结果，不展开该章的全部结构证明。
对边界使用低一维结论，并结合
\(\partial T_j\rightharpoonup\partial T\)，得到边界同样为局部整数重数可整流流。
\end{proofsketch}

将这条几何闭合定理接到已经完整证明的解析抽取命题上，即得到常用的紧性形式。

\begin{theorem}[Federer--Fleming 紧性]\label{b1:thm:app-d-federer-fleming-compactness}
设 \(K\subset\mathbb R^n\) 紧，\(T_j\in\mathbf I_k(\mathbb R^n)\)，且
\[
 \operatorname{supp}T_j\subset K,\quad
 \sup_j\bigl(\mathbf M(T_j)+\mathbf M(\partial T_j)\bigr)<\infty.
\]
则存在积分流 \(T\)，其支撑包含于 \(K\)，以及子列使
\[
 T_{j_\ell}\rightharpoonup T,\quad
 \mathbb F(T_{j_\ell}-T)\longrightarrow0.
\]
质量与边界质量满足相应的下半连续不等式。
\end{theorem}
\begin{proof}
由\cref{b1:prop:app-d-normal-weak-compactness}取得正规流弱极限，
其支撑仍包含于 \(K\)。
由\cref{b1:thm:app-d-integral-closure}，极限及其边界具有整数重数可整流表示。
全域质量与边界质量有限，故极限为积分流。
再由\cref{b1:prop:app-d-normal-weak-flat}得到平坦范数收敛。
两条下半连续不等式已在弱抽取过程中证明。
\end{proof}

首次把这条结果称为%
\term[federer-fleming-jin]{Federer--Fleming 紧性定理}[Federer--Fleming compactness theorem]%
时，应把它的两个控制量同时记住。
只有面积上界并不能限制越来越碎的边界；只有边界上界也不能限制不断增加的闭曲面。
固定紧支撑防止整个对象逃向无穷远。
在一般开集中的局部版本可以用紧集穷尽和对角抽取得到，
但局部收敛不能保证全空间的平坦范数收敛，不能省去无穷远处的控制。

% !TeX root = ../../Book1.tex
\section{填充估计与几何变分问题}
\label{b1:sec:D106}

积分流紧性最直接的用途，是从一列面积越来越小的候选曲面中取得极小者。
但在抽取子列之前，需要先回答两个问题：是否至少有一个候选对象，
以及能否把全部候选限制在某个共同紧集内。
锥构造解决第一类问题；在欧氏空间中，闭球上的不增距投影处理第二类问题。

\subsection{等周填充与尺度}
\label{b1:subsec:D10601}

锥填充的估计含有承载集的半径：
\(\mathbf M(C_aT)\leq R\mathbf M(T)/(k+1)\)。
若只知道边界的质量，不知道其形状和直径，就需要一个与半径无关的估计。
变形定理与网格整数系数的离散性恰能给出它。

\begin{theorem}[Federer--Fleming 等周不等式]\label{b1:thm:app-d-isoperimetric}
设 \(1\leq k<n\)，\(T\in\mathbf I_k(\mathbb R^n)\) 紧支撑且 \(\partial T=0\)。
则存在紧支集 \(R\in\mathbf I_{k+1}(\mathbb R^n)\)，满足
\[
 \partial R=T,\quad
 \mathbf M(R)\leq C(n,k)\,\mathbf M(T)^{(k+1)/k}.
\]
\end{theorem}
\begin{proof}
零流可取零填充，以下设 \(\mathbf M(T)>0\)。
在\cref{b1:thm:app-d-deformation}中，因 \(\partial T=0\)，
边界误差 \(S\) 的质量为零，所以 \(S=0\)，并且 \(\partial P=0\)。
记变形定理中控制 \(\mathbf M(P)\) 的常数为 \(C_0\)，选取网格尺度
\[
 \delta=\bigl(2C_0\mathbf M(T)\bigr)^{1/k}.
\]
若 \(P\neq0\)，整数网格流的最小质量性质给出
\[
 \mathbf M(P)\geq\delta^k=2C_0\mathbf M(T),
\]
与 \(\mathbf M(P)\leq C_0\mathbf M(T)\) 矛盾。
因此 \(P=0\)，变形恒等式成为 \(T=\partial R\)。
最后
\[
 \mathbf M(R)\leq C\delta\mathbf M(T)
        \leq C'\mathbf M(T)^{(k+1)/k}.
\]
变形定理的支撑控制保证该填充有紧支集。
\end{proof}

这一证明沿用 Simon 的《Lectures on Geometric Measure Theory》%
\cite{Simon1983Lectures}第30.1条的网格选择。
指数也可由缩放检验：把距离乘以 \(\lambda\)，边界质量乘以 \(\lambda^k\)，
填充质量乘以 \(\lambda^{k+1}\)，因此两侧的指数必须匹配。

整数条件在证明中用于“非零网格流至少含一整块面”。
容许任意小实系数时，这个阈值消失，上述无半径估计不能保留相同的齐次形式：
对一个固定非零边界乘以很小的实数，填充的最小质量按一次齐次变化，
而右端按大于一的幂次变化。
零维边界也须单独处理。这时 \(\delta_b-\delta_a\) 的质量恒为 \(2\)，
连接它的线段长度却可以任意大，因此不可能只用零维边界质量控制填充长度。

\subsection{Lipschitz 推前及闭球投影}
\label{b1:subsec:D10602}

\secref{b1:sec:D103}只定义了光滑映射的推前。
现在已有积分流紧性，可以用光滑逼近把推前延拓到 Lipschitz 映射，
同时保留整数结构。这个次序避免在不光滑映射下直接对一般分布作拉回。

\begin{proposition}\label{b1:prop:app-d-lipschitz-pushforward}
设 \(T\in\mathbf I_k(\mathbb R^n)\) 有紧支集，
\(f:\mathbb R^n\rightarrow\mathbb R^N\) 为 Lipschitz 映射。
则存在唯一由光滑等 Lipschitz 逼近所确定的积分流 \(f_\#T\)，并满足
\[
 \partial(f_\#T)=f_\#(\partial T),\quad
 \operatorname{supp}(f_\#T)\subset f(\operatorname{supp}T),
\]
以及
\[
 \mathbf M(f_\#T)\leq\operatorname{Lip}(f)^k\,\mathbf M(T).
\]
对 \(k=0\)，右侧的 Lipschitz 幂因子约定为 \(1\)。
推前是线性的；若 \(f=g\) 在 \(\operatorname{supp}T\) 上成立，则 \(f_\#T=g_\#T\)。
\end{proposition}
\begin{proof}
先给出逼近独立性所需的估计。
若 \(f_1,f_2\) 光滑，Lipschitz 常数不超过 \(L\)，
令 \(A=\max\{1,L\}\)，并在
\(K=\operatorname{supp}T\) 上记
\(\varepsilon=\sup_K|f_1-f_2|\)。
取直线同伦 \(h(t,x)=(1-t)f_1(x)+tf_2(x)\)。
其时间导数在 \(K\) 上不超过 \(\varepsilon\)，空间导数不超过 \(L\)。
同伦公式及质量估计给出
\[
 \mathbb F\bigl((f_2)_\#T-(f_1)_\#T\bigr)
 \leq \varepsilon A^k\bigl(\mathbf M(T)+\mathbf M(\partial T)\bigr).
\]
零维时边界项为零，同一个上界仍有效。

对 \(f\) 逐分量磨光，得到 \(f_j\rightarrow f\) 一致、
\(\operatorname{Lip}(f_j)\leq\operatorname{Lip}(f)\) 的光滑映射。
光滑推前 \((f_j)_\#T\) 的质量及边界质量一致有界，
其支撑都落在 \(f(K)\) 的一个固定有界闭邻域内。
由积分流紧性定理，存在弱收敛子列，其极限为积分流。
上述同伦估计使全部逼近项构成平坦范数 Cauchy 列，
所以任意弱收敛子列的极限都相同。
再与一个给定收敛子列比较，可知整列平坦收敛于该极限。
这里没有预先假定整个平坦空间完备，而是由紧性提供实际存在的极限。

任意另一个光滑逼近族，只要在 \(K\) 上一致逼近且 Lipschitz 常数一致受控，
用同伦估计交叉比较便得到同一个极限。
若 \(f=g\) 于 \(K\)，其两个逼近族在 \(K\) 上的差趋零，也给出同一推前。
线性由对相同逼近族逐项取极限得到。
边界交换式由光滑情形及边界的弱连续性得到；
支撑包含由一致逼近和弱极限的支撑性质得到。
最后使用光滑推前的质量估计及下半连续性，保留精确常数
\(\operatorname{Lip}(f)^k\)。
\end{proof}

设 \(B=\overline{B(a,R)}\)。到 \(B\) 的最近点投影为
\[
 p(x)=a+\min\!\left\{1,\frac{R}{|x-a|}\right\}(x-a),
\]
在 \(x=a\) 处取 \(p(a)=a\)。
它在球内等于恒等映射，且 \(\operatorname{Lip}(p)\leq1\)。
可以不用微分来核对后一性质：
凸集投影满足 \((x-p(x))\cdot(z-p(x))\leq0\) 对所有 \(z\in B\) 成立；
分别代入 \(z=p(y)\) 与 \(z=p(x)\)，相加得
\[
 |p(x)-p(y)|^2\leq (x-y)\cdot(p(x)-p(y))
                \leq |x-y|\cdot|p(x)-p(y)|.
\]
因此把积分流投到闭球不会增加质量。
只要其边界已经位于球内，边界推前仍等于原边界。
投影可能在球面处不可微，但这已经由刚才的延拓处理，不需要假设它光滑。

\subsection{Plateau 问题的存在性}
\label{b1:subsec:D10603}

给定一个边界，寻找质量最小的积分流，是%
\term[plateau-wenti]{Plateau 问题}[Plateau problem]%
的一种精确表述。
它保留取向与整数重数，不等同于在所有闭集合中直接最小化 Hausdorff 测度，
也不把候选对象限制为嵌入光滑曲面。

\begin{theorem}[Plateau 问题的积分流形式]\label{b1:thm:app-d-plateau-existence}
设 \(1\leq k\leq n\)，\(S\in\mathbf I_{k-1}(\mathbb R^n)\) 有紧支集且
\(\partial S=0\)。
若 \(k=1\)，另要求零维流 \(S\) 的全部整数系数之和为零。
则存在紧支集积分流 \(T_*\in\mathbf I_k(\mathbb R^n)\)，使
\[
 \partial T_*=S,\quad
 \mathbf M(T_*)=
   \inf\{\,\mathbf M(T):T\in\mathbf I_k(\mathbb R^n),\
                      \operatorname{supp}T\text{ 紧},\ \partial T=S\,\}.
\]
若 \(S\) 的支撑包含在某闭球内，可以选取 \(T_*\) 的支撑也包含在该球内。
\end{theorem}
\begin{proof}
取包含 \(\operatorname{supp}S\) 的闭球 \(B=\overline{B(a,R)}\)。
当 \(k\geq2\) 时，锥 \(C_aS\) 具有边界 \(S\)，且支撑位于 \(B\)。
当 \(k=1\) 时，写 \(S=\sum_iq_i\delta_{x_i}\)，
由总系数条件，
\[
 T_0=\sum_iq_i[a,x_i],\quad
 \partial T_0=\sum_iq_i(\delta_{x_i}-\delta_a)=S.
\]
这也是球内的有限质量积分流。因此容许类非空，质量下确界 \(m\) 有限。

选取紧支集积分流 \(T_j\)，满足 \(\partial T_j=S\) 且
\(\mathbf M(T_j)\rightarrow m\)。
这些流未必有共同紧支撑，但上一小节的投影流 \(p_\#T_j\) 全部支撑于 \(B\)，
边界仍为 \(S\)，质量不增。
因 \(m\) 是原容许类的下确界，其质量也趋于 \(m\)。
这就把最小化列置于同一个紧集内。

积分流紧性给出子列弱收敛于积分流 \(T_*\)，支撑位于 \(B\)。
边界的弱连续性给出 \(\partial T_*=S\)。
由质量下半连续性，
\[
 \mathbf M(T_*)\leq\liminf_j\mathbf M(p_\#T_j)=m.
\]
而 \(T_*\) 属于容许类，所以其质量又不小于 \(m\)，从而取到最小值。
\end{proof}

零维边界的系数和条件不仅足够，而且必要。
若 \(S=\partial T\) 且 \(T\) 紧支撑，取在该支撑附近恒为一的测试函数 \(\phi\)，则
\[
 S(\phi)=T(d\phi)=0.
\]
左端就是 \(S\) 的总系数。
一个单独点质量没有紧支集填充；两个系数相反的点质量则有有向线段填充。

存在性证明没有说明极小者唯一，也没有说明其支撑处处光滑。
高余维与多重数可能产生真实奇性，边界的几何性质又会带来额外问题。
Simon 的《Lectures on Geometric Measure Theory》%
\cite{Simon1983Lectures}第33节以后转向极小流的局部性质；
阅读这些内容之前，应先分清“最小化列存在极限”和“极限具有何种正则性”。

\subsection{闭形式给出的质量下界}
\label{b1:subsec:D10604}

紧性证明保证某个极小者存在，但通常不能直接写出它。
有时，一个闭微分形式已经给出足够尖锐的下界，可以识别具体极小者。

\begin{proposition}\label{b1:prop:app-d-closed-form-minimality}
设 \(T_0\in\mathbf I_k(\mathbb R^n)\) 紧支撑。
若存在光滑闭 \(k\) 次形式 \(\omega=d\eta\)，使
\[
 \sup_x\|\omega(x)\|_{\mathrm{co}}\leq1,\quad
 T_0(\omega)=\mathbf M(T_0),
\]
则对每个紧支集积分流 \(T\) 满足 \(\partial T=\partial T_0\)，都有
\(\mathbf M(T)\geq\mathbf M(T_0)\)。
这里不要求 \(\omega,\eta\) 本身紧支撑，配对时在各流的紧支撑附近取截断。
\end{proposition}
\begin{proof}
对紧支集闭流 \(T-T_0\)，取在其支撑邻域恒为一的光滑紧支集函数 \(\chi\)。
因 \(d\chi=0\) 在该邻域，
\[
 (T-T_0)(\omega)
    =(T-T_0)\bigl(d(\chi\eta)\bigr)
    =\partial(T-T_0)(\chi\eta)=0.
\]
于是 \(T(\omega)=T_0(\omega)=\mathbf M(T_0)\)。
余质量上界又给出 \(T(\omega)\leq\mathbf M(T)\)，得到结论。
计算 \(T(\omega)\) 时可选取 \(0\leq\chi\leq1\) 的截断，
因此不会破坏质量上界。
\end{proof}

虽然命题称作闭形式下界，这里还明确要求 \(\omega=d\eta\)，
以便不借助额外的上同调理论就说明所有竞争者积分相同。
在全空间上可用 Poincaré 引理从闭性推出这一点；
在一般区域中则不能省去相应条件。

例如令 \(D\) 是坐标 \(k\) 平面内的一个有界光滑区域，取标准定向，
\[
 \omega=dx^1\wedge\cdots\wedge dx^k
       =d\bigl(x^1\,dx^2\wedge\cdots\wedge dx^k\bigr).
\]
当 \(k=1\) 时，右侧括号理解为零次形式 \(x^1\)。
这个形式的余质量为一，并在 \(D\) 的取向上恒取值一，
故 \([D]\) 在所有具有相同边界的紧支集积分流中质量最小。
特别地，平面圆盘是其边界圆的一张最小质量积分流。
这项论证比较的是所有积分流竞争者，不只是图形曲面或光滑小扰动。

最后也可把固定边界换成惩罚项。
设 \(K\) 紧，\(S\in\mathbf I_{k-1}(\mathbb R^n)\) 支撑于 \(K\)，
对固定 \(\lambda>0\) 考虑
\[
 \mathcal E_\lambda(T)=\mathbf M(T)+\lambda\mathbf M(\partial T-S),
 \quad T\in\mathbf I_k(\mathbb R^n),\
                 \operatorname{supp}T\subset K.
\]
零流使容许类非空，且能量有界控制
\[
 \mathbf M(T)+\mathbf M(\partial T)
 \leq (1+\lambda^{-1})\mathcal E_\lambda(T)+\mathbf M(S).
\]
于是最小化列满足积分流紧性；
边界误差的弱收敛及质量下半连续性保证能量也取到最小值。
这一模型允许不完全匹配边界，但必须为误差支付质量代价。
它与精确 Plateau 问题的容许类不同，不能由有限的 \(\lambda\) 自动断言极小者满足
\(\partial T=S\)。

从 BV 到流，变分法的解析骨架仍是有界性、紧性和下半连续性。
新增的几何困难是：极限能否保留维数、取向和整数重数。
切片、变形及整数闭合定理承担了这部分工作，而不是把它们交给一个没有几何信息的弱极限。


\begin{chaptersummary}
切测度描述局部放大后的质量分布，流则进一步保存取向，并用测试形式上的外微分定义边界。
余质量及其对偶质量把有符号的积分恢复为无符号的大小；对可整流流，这个大小就是重数加权面积。
有限周长集合的 Gauss--Green 公式因此成为流的边界表示。

同伦与切片分别记录移动和切割产生的误差。
平坦范数允许用高一维填充计量这些误差，使靠近的曲面可以趋近，
而不要求它们的承载集逐点相同。
质量和边界质量提供正规流的解析紧性；
变形及整数闭合定理补充保留整数几何结构所需的部分。
由此得到的极小流是变分问题的解，但其光滑性、奇点和唯一性仍是进一步研究的问题。
\end{chaptersummary}

\BookExercises

% 自拟。目标：从混合测度中辨认原子主尺度；区别一般切测度系数与固定幂归一化。
\begin{exercise}\label{b1:ex:app-d-atomic-tangents}
设 \(\mu=a\delta_0+\nu\)，其中 \(a>0\)，\(\nu\) 是
\(\mathbb R^n\) 上的正 Radon 测度且 \(\nu(\{0\})=0\)。
证明 \(\operatorname{Tan}(\mu,0)=\{\,b\delta_0:b>0\,\}\)。
若 \(\nu=\mathcal H^k\llcorner P\)，其中 \(k\geq1\)、\(P\) 为过原点的 \(k\) 维平面，
说明为什么 \(r^{-k}\) 归一化不能产生非零 Radon 淡极限。
\end{exercise}
\begin{solution}
对支撑于 \(B(0,R)\) 的连续测试函数 \(\varphi\)，
\[
 \left|\int\varphi(y/r)\,d\nu(y)\right|
       \leq\|\varphi\|_\infty\,\nu(B(0,Rr))\longrightarrow0.
\]
最后一步来自 \(\nu\) 在某固定球上有限及向下连续性。
所以 \((\eta_{0,r})_\#\mu\rightarrow a\delta_0\)，乘以任意正常数可得到所需射线。

反之，设 \(c_j(\eta_{0,r_j})_\#\mu\rightarrow\lambda\neq0\)。
取非负紧支集函数 \(\varphi\)，使 \(\varphi(0)=1\)。
被测试的积分至少为 \(ac_j\)，故 \(c_j\) 必有界。
若有趋零子列，由前面的局部质量计算，沿该子列所有紧支集测试积分都趋零，
与 \(\lambda\neq0\) 矛盾。
再抽子列令 \(c_j\rightarrow c>0\)，上述 \(\nu\) 部分的估计给出
\(\lambda=ca\delta_0\)。
这就排除了任何非点质量切测度。

在后一情形，固定 \(\varphi(0)=1\) 的非负测试函数满足
\[
 r^{-k}\int\varphi(y/r)\,d\mu(y)\geq ar^{-k}\longrightarrow\infty.
\]
所以它不可能淡收敛于局部有限测度；同一点上原子项比 \(k\) 维面积项更强。
\end{solution}

% 自拟。目标：构造非测度型零维流，并发现有限平坦范数不蕴含有限质量。
\begin{exercise}\label{b1:ex:app-d-dirac-derivative}
在 \(\mathbb R\) 上令 \(T(\varphi)=\varphi'(0)\)。
证明 \(T\) 是零维流，\(\mathbf M(T)=\infty\)，但 \(\mathbb F(T)\leq1\)。
这是否与零维积分流是有限整数点质量和的描述矛盾？
\end{exercise}
\begin{solution}
在任何共同紧支集上，\(|T(\varphi)|\leq\|\varphi'\|_\infty\)，
所以 \(T\) 是连续线性泛函。
选取 \(\psi\in C_c^\infty(\mathbb R)\)，满足
\(\|\psi\|_\infty\leq1\)、\(\psi'(0)\neq0\)。
对 \(\varphi_j(x)=\psi(jx)\)，有 \(\|\varphi_j\|_\infty\leq1\)，
而 \(|T(\varphi_j)|=j|\psi'(0)|\rightarrow\infty\)，故质量无限。

定义一维流 \(B(g\,dx)=g(0)\)。其质量为 \(1\)，并且
\(\partial B(\varphi)=B(d\varphi)=\varphi'(0)=T(\varphi)\)。
取平坦分解 \(T=0+\partial B\)，便有 \(\mathbb F(T)\leq1\)。
这个 \(T\) 不是积分流，甚至不是有限质量零维流；
而 \(B\) 也不正规，因为其边界质量无限。
“零维流”“有限质量零维流”“零维积分流”是逐步缩小的类别，没有矛盾。
\end{solution}

% 自拟。目标：将分布的零导数性质用于最高维流，得到固定边界的唯一性特例。
\begin{exercise}\label{b1:ex:app-d-top-dimensional-uniqueness}
证明：若 \(T\) 是 \(\mathbb R^n\) 上有限质量的 \(n\) 维流且 \(\partial T=0\)，
则 \(T=0\)。
推出具有相同边界的两个有限质量 \(n\) 维流必相同。
这里为什么不能把 \(n\) 维改成任意 \(k<n\)？
\end{exercise}
\begin{solution}
最高次形式只有一个系数，因此 \(T\) 对应一个有限符号测度 \(\mu\)：
\[
 T(\varphi\,dx^1\wedge\cdots\wedge dx^n)=\int\varphi\,d\mu.
\]
对省去 \(dx^i\) 的测试 \(n-1\) 次形式使用 \(\partial T=0\)，
得到 \(\int\partial_i\varphi\,d\mu=0\)，即 \(\partial_i\mu=0\) 作为分布成立。
用光滑核卷积后，函数 \(u_\varepsilon=\rho_\varepsilon*\mu\)
的全部偏导都为零，所以它在连通的 \(\mathbb R^n\) 上为常数。
又有
\(\|u_\varepsilon\|_{L^1}\leq|\mu|(\mathbb R^n)<\infty\)，
全空间上的可积常数只能为零。
令 \(\varepsilon\downarrow0\)，在测试函数上恢复 \(\mu=0\)，故 \(T=0\)。

对两个具有相同边界的流取差即可得到唯一性。
低维时，一条有向圆周给出非零、有限质量、零边界的一维流；
更一般的球面边界也提供反例。
其质量集中在低维几何对象上，不对应全空间常数密度。
\end{solution}

% 自拟。目标：构造质量下半连续而不连续的几何例子，使用高一维薄层填充。
\begin{exercise}\label{b1:ex:app-d-thin-rectangles}
令 \(Q_j=(0,1)\times(0,j^{-1})\subset\mathbb R^2\)，
\(T_j=\partial[Q_j]\)。
求 \(\mathbf M(T_j)\) 与 \(\partial T_j\)，证明 \(T_j\) 按平坦范数趋于零，
并解释该例与质量下半连续性的关系。
\end{exercise}
\begin{solution}
矩形边界各边内部互不重合，单位重数给出
\[
 \mathbf M(T_j)=2+2j^{-1},\quad \partial T_j=0.
\]
取填充 \(B=[Q_j]\)，其质量为 \(j^{-1}\)，所以
\[
 \mathbb F(T_j)\leq j^{-1}\longrightarrow0.
\]
因此 \(T_j\rightharpoonup0\)，但其质量趋于 \(2\)，并不趋于极限流的质量 \(0\)。
下半连续性要求的是
\(\mathbf M(0)\leq\liminf_j\mathbf M(T_j)\)，在这里是严格不等式。
几何上，两条较长的边反向靠拢并相消；有向极限允许抵消，不保存无向周长总和。
\end{solution}

% 自拟。目标：说明积分流紧性中的边界质量界不能删除；极限出现半整数重数。
\begin{exercise}\label{b1:ex:app-d-unbounded-boundary}
把 \([0,1]\) 嵌入 \(\mathbb R^2\) 的横轴，令
\[
 T_j=\sum_{i=0}^{j-1}
       \left[\left(\frac{i}{j},0\right),
             \left(\frac{i+1/2}{j},0\right)\right].
\]
证明 \(T_j\) 的质量和支撑一致受控，
但它弱收敛于重数为 \(1/2\) 的整条单位线段流。
计算边界质量，并说明哪个紧性假设失败。
\end{exercise}
\begin{solution}
各短线段内部互不相交且取向相同，故
\[
 \mathbf M(T_j)=\frac12,\quad
 \operatorname{supp}T_j\subset[0,1]\times\{0\}.
\]
每条线段有一个正端点和一个负端点，所有 \(2j\) 个端点互不重合，
所以 \(\mathbf M(\partial T_j)=2j\)。

对测试一形式 \(\omega=a(x,y)\,dx+b(x,y)\,dy\)，只有第一项贡献积分。
在每个区间 \([i/j,(i+1)/j]\) 内，以 \(a(i/j,0)\) 代替 \(a(x,0)\)，
两种积分的误差都由该系数在尺度 \(1/j\) 上的连续模控制。
所有区间相加后得到
\[
 T_j(\omega)-\frac12\int_0^1a(x,0)\,dx\longrightarrow0.
\]
故弱极限为 \(\frac12[(0,0),(1,0)]\)。
其正长度部分的重数为 \(1/2\)，不是整数，因而不是积分流。
失败的是边界质量的一致有界性；仅凭原流面积界不能在极限中保留整数层数。
\end{solution}

% 自拟。目标：区分局部弱收敛与全空间平坦收敛，检验无穷远紧性。
\begin{exercise}\label{b1:ex:app-d-escape-to-infinity}
设 \(T\neq0\) 是 \(\mathbb R^2\) 上沿单位圆周的积分流，
\(T_j=(\tau_{(3j,0)})_\#T\)。
证明 \(T_j\rightharpoonup0\)，但没有平坦范数收敛子列。
这些流的质量与边界质量是否仍一致有界？
\end{exercise}
\begin{solution}
任意固定测试形式具有紧支集，充分大的 \(j\) 使它与平移圆周不相交，
因此 \(T_j\) 在该测试形式上最终为零。
另一方面，
\(\mathbf M(T_j)=2\pi\)、\(\partial T_j=0\)，两个质量仍一致有界。

欧氏平移保持质量和边界，并将每个平坦分解推前为同样代价的分解。
对逆平移再作一次比较，得到
\(\mathbb F(T_j)=\mathbb F(T)>0\)。
若有子列按平坦范数收敛，其弱极限必须是零；
但按范数收敛到零要求上述范数趋零，矛盾。
因此共同紧支撑不能仅用各个支撑分别紧来代替。
\end{solution}

% 自拟。目标：计算零维平坦距离，区分移动代价与直接删除代价。
\begin{exercise}\label{b1:ex:app-d-flat-two-points}
对不同的 \(a,b\in\mathbb R^n\)，证明
\[
 \mathbb F(\delta_b-\delta_a)=\min\{\,|b-a|,2\,\}.
\]
由此说明平坦范数在空间伸缩下为什么不总按单一维数幂次变化。
\end{exercise}
\begin{solution}
把整个差作为质量项，可得上界 \(2\)；
把它写成有向线段的边界，可得上界 \(|b-a|\)。

记 \(d=|b-a|\)、\(e=(b-a)/d\)，考虑一维截断函数
\[
 q(s)=\max\{-1,\min\{1,s-d/2\}\}.
\]
函数 \(q(e\cdot(x-a))\) 的绝对值不超过一、Lipschitz 常数不超过一，
且在 \(b,a\) 两点之差恰为 \(\min\{d,2\}\)。
为得到紧支集测试函数，先乘一个在包含这两点的大球内为一、
梯度上界趋零的光滑截断；其 Lipschitz 常数至多为 \(1+\varepsilon\)。
再磨光并除以 \(1+\varepsilon\)。
这样得到 \(\varphi_\varepsilon\in C_c^\infty\)，
\[
 \|\varphi_\varepsilon\|_\infty\leq1,\quad
 \|\nabla\varphi_\varepsilon\|_\infty\leq1,
\]
且两点之差趋于 \(\min\{d,2\}\)。
平坦范数的测试估计因此给出反向下界。

把两点距离乘以 \(\lambda\)，所得距离为 \(\min\{\lambda d,2\}\)。
短距离时选择线段填充，长距离时直接支付两个点的质量更便宜；
分解中两个相邻维数的代价竞争，使它不具有单一的空间缩放幂律。
\end{solution}

% 自拟。目标：利用整数零维误差的离散性求解一个边界惩罚问题，联系距离下界。
\begin{exercise}\label{b1:ex:app-d-penalty-threshold}
设 \(a\neq b\)，紧集 \(K\subset\mathbb R^n\) 包含线段 \([a,b]\)，\(\lambda>0\)。
在支撑于 \(K\) 的积分一维流上，证明
\[
 \min_T\left\{\mathbf M(T)
       +\lambda\mathbf M\bigl(\partial T-(\delta_b-\delta_a)\bigr)\right\}
       =\min\{\,|b-a|,2\lambda\,\}.
\]
说明当 \(2\lambda<|b-a|\) 时零流为何是唯一极小者，
而当 \(2\lambda>|b-a|\) 时所有极小者都精确满足指定边界。
\end{exercise}
\begin{solution}
设误差流 \(R=\partial T-(\delta_b-\delta_a)\)。
它是有限整数系数零维流，且全部系数之和为零：
对 \(\partial T\) 用在 \(K\) 邻域恒为一的测试函数计算即可。
若 \(R\neq0\)，正系数和与负系数和的绝对值至少各为一，
因此 \(\mathbf M(R)\geq2\)，能量至少为 \(2\lambda+\mathbf M(T)\)。

若 \(R=0\)，令 \(e=(b-a)/|b-a|\)。
常值一形式 \(\omega=e\cdot dx=d(e\cdot x)\) 的余质量为一。
在 \(T\) 的紧支撑附近取截断，边界公式给出
\[
 T(\omega)=(\delta_b-\delta_a)(e\cdot x)=|b-a|.
\]
所以 \(\mathbf M(T)\geq|b-a|\)。
这两种情形共同给出所需下界；零流的能量为 \(2\lambda\)，
有向线段流的能量为 \(|b-a|\)，下界确能达到。

若 \(2\lambda<|b-a|\)，精确边界的候选不可能极小，
而非零误差候选要达到 \(2\lambda\)，必须满足 \(\mathbf M(T)=0\)，即 \(T=0\)。
若 \(2\lambda>|b-a|\)，任何非零误差都已经支付过高代价，极小者只能满足 \(R=0\)。
在相等的临界值，零流与线段流都为极小者；
有限惩罚参数并不无条件等同于精确边界约束。
\end{solution}
